\section{组合结构}

基准稳健核心版按资产层而不是按单只股票构建。组合内受保障流动性与交易现金合计25\%，用来降低被动卖出风险；利率债为35\%，红利权益ETF为30\%，非杠杆黄金ETF为10\%。直接股票和公募REITs在初始向量中均为0\%。目标权重不等于当前已通过实施状态：截至本报告，ETF资格包未运行，未通过额度按规则留在受保障流动性。

\begin{exhibitbox}[基准稳健核心版]
\small
\begin{tabularx}{\exhibitboxwidth}{L{3.1cm}P P P}
\toprule
\textbf{资产层} & \textbf{目标权重} & \textbf{固定金额} & \textbf{资格与上限} \\
\midrule
组合内受保障流动性 & 20\% & CNY40,000 & 合格存款或到期前直接政府债；唯一承接未通过ETF资格检查的余额，安全上限100\%。 \\
交易现金 & 5\% & CNY10,000 & 硬上限CNY10,000；货币工具不自动等同于受保障存款，也不承接延后ETF资金。 \\
高质量人民币利率债 & 35\% & CNY70,000 & 指定\texttt{511010}；产品、久期和交易质量$q=1$后才可实施。 \\
红利权益ETF A & 15\% & CNY30,000 & 指定\texttt{510880}；产品与交易质量$q=1$后才可分批。 \\
红利权益ETF B & 15\% & CNY30,000 & 指定\texttt{512890}；另须通过合并行业和集团穿透，否则保持0。 \\
非杠杆黄金ETF & 10\% & CNY20,000 & 指定\texttt{518880}；非杠杆、价格/IOPV和交易质量$q=1$后才可分批，硬上限12\%。 \\
公募REITs & 0\% & CNY0 & 仅在三只独立项目同时通过时，从利率债转入最多5\%。 \\
直接高股息股票 & 0\% & CNY0 & 观察池；未来只能替换红利权益ETF，合计上限15\%。 \\
\midrule
\textbf{合计} & \textbf{100\%} & \textbf{CNY200,000} & \\
\bottomrule
\end{tabularx}
\end{exhibitbox}

\section{指定模型标的与资格路径}

执行日资格包是一个可复核的字段清单，而非概括性的“交易质量”标签。每只ETF需记录同一中国市场时段的产品文件、来源、时间戳、价格和IOPV/参考值、折溢价、最佳买卖价、价差、成交额或可执行深度、最小交易单位和账户可行性。模型控制要求绝对折溢价不高于1.00\%、买卖价差不高于中间价的0.50\%、单笔不高于记录的同日成交额或可执行深度的10\%。这些是风险控制规则，不是当前市场观察。

\begin{exhibitbox}[指定模型标的]
\small
\begin{tabularx}{\exhibitboxwidth}{L{2.4cm}L{3.1cm}L{1.7cm}P}
\toprule
\textbf{资产层} & \textbf{指定标的} & \textbf{目标金额} & \textbf{资格条件与失败时动作} \\
\midrule
高质量人民币利率债 & \texttt{511010} 国泰上证5年期国债ETF & CNY70,000 & 当前产品文件、久期、价格/IOPV、折溢价、价差和流动性均通过使$q=1$；否则不交易，CNY70,000留在受保障流动性。 \\
红利权益ETF A & \texttt{510880} 华泰柏瑞上证红利ETF & CNY30,000 & 当前产品文件、价格/IOPV、折溢价、价差和流动性均通过使$q=1$；失败批次留在受保障流动性。 \\
红利权益ETF B & \texttt{512890} 华泰柏瑞中证红利低波动ETF & CNY30,000 & 除交易资格外，合并行业不得超过总组合25\%，前十大集团不得超过权益桶35\%；不通过则保持0，CNY30,000留在受保障流动性。 \\
非杠杆黄金ETF & \texttt{518880} 华安易富黄金ETF & CNY20,000 & 核验非杠杆现货暴露、文件、IOPV、折溢价、价差和流动性；失败批次留在受保障流动性。 \\
交易现金工具 & 账户结算现金或银行活期；\texttt{511880}只作可选补充 & CNY10,000 & 交易现金上限为CNY10,000；若使用\texttt{511880}，另核验申赎/交易机制、价格偏离和非存款风险。 \\
\bottomrule
\end{tabularx}
\end{exhibitbox}

\section{六个月建仓动作}

每个金额是对CNY200,000的目标名义金额；场内最小交易单位和当日价格导致的零头向下取整，并留在受保障流动性。任何延后金额不为凑足目标而增加交易现金或换入相关资产。当前资格包未运行，故下表为通过资格后才会进入的计划顺序，而不是现在的交易指令。

\begin{exhibitbox}[六个月模型建仓动作]
\small
\begin{tabularx}{\exhibitboxwidth}{L{2.5cm}P P}
\toprule
\textbf{阶段} & \textbf{金额与目标} & \textbf{前提与失败处理} \\
\midrule
T0防守层 & CNY40,000受保障流动性、CNY10,000交易现金、CNY70,000的\texttt{511010} & 前两项按账户和存款边界核验；\texttt{511010}仅在$q=1$时建立，未通过的CNY70,000留在受保障流动性。 \\
第1--5个月 & 每月CNY5,000的\texttt{510880}、CNY5,000的合格\texttt{512890}、CNY3,333的\texttt{518880} & 各ETF分别$q=1$；\texttt{512890}还须穿透合格。任一失败时仅延后该笔，并将该笔留在受保障流动性。 \\
第6个月 & CNY5,000的\texttt{510880}、CNY5,000的合格\texttt{512890}、CNY3,335的\texttt{518880} & 以可交易最小单位向下取整；所有余量与未通过金额均留在受保障流动性，不产生补单。 \\
公募REITs或直接股票 & CNY0 & 不因等待分配、市场下跌或风险资产延后自动启用；须独立通过预先设定门槛。 \\
\bottomrule
\end{tabularx}
\end{exhibitbox}

\section{追加与再平衡}

每季度记录实际到账现金$D_0$、资产权重、产品事件、集中度、外部现金流和实际回撤$DD_t$；每半年才进行常规再平衡，优先使用新增资金和实际分配，避免把价格下跌自动解释为加仓信号。以下情形启动人工复核：受保障流动性、利率债或总权益相对目标偏离超过5个百分点；交易现金超过CNY10,000或在分批期内低于目标；黄金偏离超过2个百分点；任何资产超过硬上限或失去资格；或有完整账户记录的$DD_t\leq-12\%$。压力情景数值不代替实际$DD_t$。
